\documentclass[11pt,a4paper]{article}
\usepackage[a4paper,left=28mm,right=38mm,top=25mm,bottom=25mm]{geometry}
\usepackage{amsmath,amssymb,amsthm,mathtools}
\usepackage{booktabs,tabularx,array,enumitem}
\usepackage[T1]{fontenc}
\usepackage{lmodern}
\usepackage{microtype}
\usepackage{xcolor}
\usepackage{setspace}
\usepackage{fancyhdr}
\usepackage[hidelinks]{hyperref}
\usepackage{xurl}
\setstretch{1.18}
\setlength{\parskip}{4pt}
\setlength{\parindent}{0pt}
\definecolor{boxgray}{gray}{0.93}
\newcommand{\statusbox}[1]{\par\smallskip\noindent\colorbox{boxgray}{\parbox{\dimexpr\linewidth-2\fboxsep}{\small #1}}\par\smallskip}
\theoremstyle{plain}
\newtheorem{claim}{Claim}
\newtheorem*{corollary}{Corollary}
\theoremstyle{definition}
\newtheorem*{setting}{Setting}
\theoremstyle{remark}
\newtheorem*{remark}{Remark}
\newcommand{\eps}{\epsilon}
\newcommand{\R}{\mathbb{R}}
\newcommand{\simp}{\Delta^{\circ}}
\pagestyle{fancy}
\fancyhf{}
\lhead{\small sense-correct-abstain}
\rhead{\small S1$''$ --- draft for owner review}
\cfoot{\small \thepage}

\begin{document}
\begin{center}
{\LARGE\bfseries Extension S1$''$}\\[4pt]
{\large Per-context shared response and visibility/offset readings of spec \S5.3}\\[8pt]
{\small Source: \texttt{docs/errata/E007\_owner\_review\_lemmas.md} \S3.2 (local commit \texttt{630ef91}); this sheet adds the exceptional-point and enlargement proofs in full. 2026-11-21.}
\end{center}

\statusbox{\textbf{Status.} Proved (draft, Claude); \emph{owner review pending}. Notation and the
classes N-loc, N-free follow \texttt{S1prime\_B\_rigorous.pdf} (owner note). Scope: synthetic
compiled-angle model, frozen Gate 1A designs, true points with $\phi_0=0$ and
$|\eps_0|\le 0.04$. Whether the readings below exhaust spec \S5.3 is a specification
question and is \emph{not} proved here.}

\section{Why an extension is needed}
Spec \S5.3: ``treats \emph{each target context's} ideal outcome distribution and its local
response directions as unknown \dots\ Visibility and label-independent offset are also nuisance
quantities.'' The owner note and E006 parameterize the nuisance \emph{per candidate}. Two
readings fall outside that:
\begin{enumerate}[label=(\alph*),itemsep=2pt]
\item \textbf{per-context shared response}: in the \texttt{distinct} design the two candidates
belong to one context and could share one response map;
\item \textbf{visibility and offset}: laws $vP+(1-v)o$.
\end{enumerate}
In the \texttt{single} and \texttt{same} designs every candidate has the same angle, so a shared
response is evaluated at one point only and the per-candidate results apply unchanged.

\section{Shared response (distinct design)}

\begin{setting}
One context, candidates $k=1,2$. Candidate $k$ uses one tagged angle $\theta_k$ for all its
$L$ tagged blocks, with $\theta_1\neq\theta_2$; write $\Theta_k:=\theta_k\mathbf 1_L$. A single
map $R:\R^L\to\simp$ from a class $\mathcal F$ serves both candidates. By frame covariance
(owner note, Lemma 2), at $\phi=0$:
\[
P_{k,t}=R\big((1+\eps)\Theta_k\big)=:P_k\qquad\text{for every label pattern } t,\ k=1,2 .
\]
Observational equivalence means equality of $(P_1,P_2)$ (labels and candidate mix are known).
Assume $1+\eps_0\neq0$ (true on the grid).
\end{setting}

\begin{claim}[S1$''$(a), nonlinear shared class]\label{c:a}
Suppose $\mathcal F$ is closed under rescaling of the argument, $R\mapsto R(c\,\cdot)$ for
every $c\neq0$ (true for all maps, continuous maps, and $C^k$ maps). Then every
$\eps'\neq-1$ belongs to the identified set $I$.
\end{claim}
\begin{proof}
\begin{enumerate}[label=\textbf{Step \arabic*.},leftmargin=*,itemsep=2pt]
\item Fix $\eps'\neq-1$ and put $c:=(1+\eps_0)/(1+\eps')$, so $c\neq0$. Define
$R'(u):=R(cu)$. By closure $R'\in\mathcal F$, and $R'$ has the regularity of $R$.
\item For $k=1,2$:
$R'\big((1+\eps')\Theta_k\big)=R\big((1+\eps_0)\Theta_k\big)=P_k$.
\item With $\phi'=0$, the point $(\eps',R')$ reproduces $(P_1,P_2)$. \qedhere
\end{enumerate}
\end{proof}

\begin{claim}[S1$''$(b), the exceptional point $\eps'=-1$]\label{c:b}
Under the setting, $-1\in I$ implies $P_1=P_2$. Conversely, if $P_1=P_2$ and $\mathcal F$
contains the constant maps, then $-1\in I$. Hence, with constants in $\mathcal F$,
$I=\R$ if $P_1=P_2$ and $I=\R\setminus\{-1\}$ if $P_1\neq P_2$.
\end{claim}
\begin{proof}
At $\eps'=-1$ both candidates are evaluated at $u=0$, so any equivalent point has
$P_1=R'(0)=P_2$. Conversely, if $P_1=P_2$, the constant map $R'\equiv P_1$ works.
Combine with Claim~\ref{c:a}.
\end{proof}

\begin{remark}
This is the reading in which the draft condition $1+\eps'\neq0$ is genuinely needed. The
owner's correction (``unnecessary'') is correct for the \emph{per-candidate} class, where the
translation $R_c(\cdot+(\eps_0-\eps')\theta_c)$ is available for each candidate separately.
With a shared map, a single translation cannot serve two candidates, because the required
shifts $(\eps_0-\eps')\theta_k$ differ.
\end{remark}

\begin{claim}[S1$''$(c), first-order shared class]\label{c:c}
Let $P_{k,t}=q+\sum_\ell\big(\eps\theta_k+s_\ell(t)\phi\big)D_\ell$ with one $q\in\simp$ and
shared, unknown $D_\ell\in\Sigma_0$ (positivity required). Then every $\eps'\neq0$ belongs
to $I$, and $0\in I$ iff $P_1=P_2$.
\end{claim}
\begin{proof}
At $\phi=0$, $P_k=q+\eps_0\theta_k d$ with $d:=\sum_\ell D_\ell\in\Sigma_0$.
For $\eps'\neq0$, set $D'_\ell:=(\eps_0/\eps')D_\ell$ and $q':=q$; then
$\eps'\theta_kd'=\eps_0\theta_kd$, so $P'_k=P_k$ for $k=1,2$, and positivity holds because
the laws are the true ones. At $\eps'=0$ every equivalent point has $P'_1=q'=P'_2$, which
requires $P_1=P_2$; conversely, if $P_1=P_2$, take $q'=P_1$ and $D'=0$.
\end{proof}

\begin{remark}
If the shared directions $D_\ell$ were \emph{known}, then
$P_1-P_2=\eps(\theta_1-\theta_2)d$ would identify $\eps$ whenever $d\neq0$. That class is not
a reading of \S5.3, which declares the response directions unknown, but it shows that the
conclusion depends on that clause.
\end{remark}

\section{Visibility and label-independent offset}

\begin{claim}[S1$''$(d), enlargement]\label{c:d}
Let $\mathcal M$ be any of the classes above, with identified set $I_{\mathcal M}(\eps_0,\eta_0)$.
Let $\mathcal M'$ have laws $vP_{k,t}+(1-v)o$ with $v\in(0,1]$ and $o\in\simp$
label- and candidate-independent (extra nuisance $(v,o)$). Then for every true $(v_0,o_0)$,
\[
I_{\mathcal M'}(\eps_0,\eta_0,v_0,o_0)\supseteq I_{\mathcal M}(\eps_0,\eta_0).
\]
\end{claim}
\begin{proof}
If $\eps'\in I_{\mathcal M}$, there is $\eta'$ with $P'_{k,t}=P_{k,t}$ for all $(k,t)$. Then
$(\eps',\eta',v_0,o_0)$ gives $v_0P'_{k,t}+(1-v_0)o_0=v_0P_{k,t}+(1-v_0)o_0$, the true
$\mathcal M'$-law. Positivity is inherited.
\end{proof}

\section{Consequence for Gate 1A}

\begin{corollary}
In each of the six primary designs (G2, G3 $\times$ single, same, distinct), at each of the nine
grid points with $\phi_0=0$, under the per-candidate classes (owner note) and under the
shared and visibility/offset readings (a)--(d), the identified set contains an open interval
around $\eps_0$. Hence the profile maximiser is not unique.
\end{corollary}
\begin{proof}
\begin{enumerate}[label=\textbf{Step \arabic*.},leftmargin=*,itemsep=2pt]
\item \texttt{single}, \texttt{same}: owner note, Theorem S1$'$(1)--(3); Claim~\ref{c:d} for
visibility/offset.
\item \texttt{distinct}, shared nonlinear: Claims~\ref{c:a}--\ref{c:b} give
$I\supseteq\R\setminus\{-1\}$, an open set containing $\eps_0\in[-0.04,0.04]$.
\item \texttt{distinct}, shared first order: Claim~\ref{c:c} gives $I\supseteq\R\setminus\{0\}$.
If $\eps_0\neq0$ this contains an open interval around $\eps_0$. If $\eps_0=0$ then
$P_1=P_2=q$ and $I=\R$.
\item Visibility/offset: Claim~\ref{c:d} preserves every inclusion.
\item By the owner note's Gibbs corollary, profile maximisers attaining $\Lambda^*$ form
exactly $I$; within the spec box $[-0.05,0.05]$, an open interval around an interior
$\eps_0$ remains a non-degenerate interval.
\end{enumerate}
\end{proof}
The final step from non-uniqueness to FAIL is the specification inference already checked
against spec \S7 and \S9 (E007 \S3.1); it is not a mathematical statement.

\section{Review checklist}
\renewcommand{\arraystretch}{1.6}
\noindent\begin{tabularx}{\linewidth}{@{}c>{\raggedright\arraybackslash}Xc@{}}
\toprule
& Item to check & Initials / date \\
\midrule
$\square$ & Readings (a)--(b) are fair interpretations of \S5.3 (spec question). & \\
$\square$ & Claim (a): rescaling construction; closure under $R\mapsto R(c\,\cdot)$, $c\neq0$. & \\
$\square$ & Claim (b): exceptional point $\eps'=-1$, iff condition $P_1=P_2$. & \\
$\square$ & Remark after (b): consistent with the owner's per-candidate correction. & \\
$\square$ & Claim (c): first-order shared version; exceptional point $\eps'=0$. & \\
$\square$ & Claim (d): enlargement by visibility/offset preserves equivalences. & \\
$\square$ & Corollary: open interval around $\eps_0$ in all six designs, including $\eps_0=0$. & \\
\bottomrule
\end{tabularx}

\vspace{14pt}
\textbf{Decision:}\quad $\square$ accept as proved\quad $\square$ accept with corrections\quad $\square$ reject

\vspace{10pt}
\textbf{Notes:}
\end{document}
